\section{The methods that actually transfer}
\label{sec:source-architecture}

The exported statements concern the Riemann zeta function and
\emph{Dirichlet} $L$-functions, including imprimitive characters. They do not
quantify over arbitrary Hecke or Artin representations; see
\src{Statements.lean}{the precise statements} and
\src{Main.lean}{the five exported theorems}~\cite{upstream}.
The relevant opportunity is to reuse the arguments inside the proof.

\subsection{One determinant, two bounds}

The source constructs integral forms whose divisibility records a weighted
prime sum, obtains many independent images by a second-moment argument,
and compares upper and lower bounds for the same Gram determinant.
Write $D$ for the target dimension, $n$ for the selected independent
vectors, $H$ for the main amount, $R$ for the prime remainder, and $m$ for
the number of lifts. With controlled errors, the comparison has the form
\[
 D(H-E_-)\le \log\det\operatorname{Gram}(y)+2\delta
       \le n(H-mR+E_+)+2\delta.
\]
Here $\delta$ measures the logarithmic index needed to make the vectors
integral. The main amount cancels, producing a bound on $R$.
The actual source conclusion is
$m\,\mathrm{totalRemainder}\le fa(\log b)^5$;
\src{New/Reduction/MainInequality.lean}{the main inequality}
identifies each contribution.

This comparison needs all three ingredients. The short forms span; their
images have controlled index; and a small set of translates spans at least
$(1-4/b^2)D$ dimensions, while its size times the source dimension is at
most $5D\log(1+2b^2)$. These are the distinct declarations
\texttt{span\_short\_integralRamified\_eq\_top},
\texttt{indexLe\_normalizedPoincare}, and \texttt{enough\_forms} in
\src{New/Reduction/ShortAndEnough.lean}{ShortAndEnough}.
An integral vector with a large divisibility exponent alone does not
provide this comparison.

The lower-bound interface is already abstract.
\texttt{TwistedLattice.covolume\_add\_index\_nonneg} in
\src{New/Reduction/Lattices/TwistedLattice.lean}{TwistedLattice}
asserts
\[
 0\le c+2\delta+n'(b_{\mathrm{short}}+b_{\mathrm{denom}})
                   +b_{\mathrm{volume}}.
\]
Here $n'$ completion vectors extend the selected family to a basis, $c$
bounds its logarithmic Gram determinant, and the remaining terms bound
completion norms, denominators, and ratios of lift determinants.
Injectivity and compatible integral lattices are hypotheses.
The theorem allows one twist \emph{per prime}, connected by a common height
estimate. Its index correction matrix must nevertheless work for every
allowed twist at that prime. The specialized lower bound is
\texttt{integral\_forms\_long} in
\src{New/Determinant/Twist/LongForms.lean}{LongForms}.

\subsection{A necessary correction concerning Selmer precision}

Some earlier transfer discussions treated the twisting result as only
providing a new character at each finite precision. The inspected source
is stronger. Its \texttt{TwistedLengthLe} predicate bounds the length of
the module with coefficients $\operatorname{Frac}(O')/O'$ itself, over
all sufficiently large coefficient rings $O'$;
see \src{New/Congruences/Selmer/Line.lean}{the definitions in Line}.
The passage from finite coefficients uses the following elementary fact.

\begin{lemma}[Finite-level stopping]
Let $O$ be a discrete valuation ring with uniformizer $\pi$, let $X$ be
a torsion $O$-module, and let $N\ge1$. If
$\operatorname{length}_O X[\pi^N]<N$, then
$X=X[\pi^N]$ and $X$ has that same finite length.
\end{lemma}
\begin{proof}
If $\pi^Nx\ne0$, choose the least $r>N$ such that $\pi^rx=0$.
The element $\pi^{r-N}x$ has exact order $\pi^N$ and generates a
length-$N$ submodule of $X[\pi^N]$, a contradiction.
\end{proof}

No finite-generation assumption on $X$ is needed. In the source, a locally
constant cocycle on the compact Galois group has finite image, so a common
denominator kills it. Finite-coefficient cocycles identify with the
corresponding torsion subgroup. Thus a strict determinant bound below the
chosen torsion precision controls the whole module
(\texttt{length\_selmerModule\_divisibleQuot\_le} in
\src{New/Congruences/Selmer/Line.lean}{Line}).
\texttt{exists\_anticyclotomic\_twist} in
\src{New/Congruences/Selmer/GeneralTheorem.lean}{GeneralTheorem}
then supplies one character whose first $m$ powers meet all the stated
conditions simultaneously. The specialized endpoint in
\src{New/Congruences/Selmer/TwistingTheorem.lean}{TwistingTheorem}
also explicitly bounds divisible coefficients.

The qualifications are arithmetic, not a missing interchange of these
precision quantifiers. A \texttt{SelmerDatum} requires residual absolute
irreducibility, a local-rank budget, finite conjugate-dual and local
defects, and an inertial element acting by a scalar $\zeta$ with
$\zeta^2-1$ invertible; see
\src{New/Congruences/Selmer/SelmerData.lean}{SelmerData}.
Its local conditions must agree with those of the intended arithmetic
group. The normalization of cocycles is legitimate when
$\rho(\gamma)-1$ is invertible: subtracting the coboundary of
$(\rho(\gamma)-1)^{-1}z(\gamma)$ gives the unique representative vanishing
at $\gamma$.

Finally, small normalized length does not mean zero. If $O'$ has absolute
ramification index $e$, the nonzero module $O'/\pi$ has normalized length
$1/e$. The twisting construction may enlarge $e$. Finiteness, vanishing,
and an exact arithmetic order therefore require different arguments.
